\chapter{Conclusion and Future Work} \label{chapter:ch7_conclusion}

\section{Conclusion}

This thesis investigated whether collaborative-robot tasks that change frequently can be programmed more directly by specifying spatial intent in the real workspace. The presented prototype suggests that the answer is positive for selected task-oriented workflows. It demonstrates that the combination of workspace sensing, spatial input, task-specific interpretation, and shared execution infrastructure can support practical task-oriented programming workflows without exposing the user to a full robot programming environment.

The main contribution is the resulting system for task-oriented spatial authoring: a modular runtime with explicit interfaces between sensing, authoring, interpretation, visualization, and execution. This system was exercised through two representative task families, pick-and-place and seam welding, and was tested in limited trials on physical hardware. Pick-and-place showed that the architecture can combine authored intent with runtime scene information, while seam welding showed more clearly where the approach is practically strongest: trajectory-centric tasks in which repeated re-teaching becomes a meaningful part of the programming cost. Together, these demonstrations suggest that spatial authoring is most promising not as a general replacement for robot programming, but for workflows in which the task changes often enough that fast re-authoring inside a prepared use case becomes valuable.

At the same time, the thesis remains a proof of concept rather than a validated industrial programming platform. The current implementation is limited by narrow scene understanding, prototype-level sensing and calibration, and an execution layer in which generalized tool control, deeper path planning, and simulation are not yet implemented. The most defensible overall conclusion is therefore not that the thesis replaces established robot-programming workflows, but that it identifies a practically interesting and technically coherent direction for task-oriented spatial authoring.

\section{Near-Term Future Work}

The clearest engineering boundary is the current scene and sensing pipeline. The first practical step is therefore to move beyond the present box-based fiducial scene model, improve calibration robustness, and make better use of the sensing already available on the pen side. It would also be useful to characterize the stability and practical accuracy of both spatial input and scene estimation more explicitly than was done in this thesis.

A second engineering boundary is the execution layer. A stronger execution layer should refine the current motion interface, design a generalized tool-control abstraction, and add execution-aware capabilities such as motion validation, collision checking, and safer planning. In the same area, it would be valuable to explore a path toward export or compilation of suitable task representations into robot-side programs, especially for use cases that do not require runtime scene adaptation.

Beyond scene, sensing, and execution, the remaining near-term engineering work is mainly completion work, such as better cross-platform support. These changes would further complete the prototype, but they do not change the main technical boundary of the system.

\section{Longer-Term Directions}

Beyond technical completion, the thesis opens several research directions. Richer interaction is the first of them. The current system uses a tracked pen, but the modular architecture makes it reasonable to explore hand tracking, VR or AR controllers, or other 6DoF inputs. An especially interesting direction is multimodal interaction that combines spatial indication with language. A system that can interpret deictic gesture together with task-specific spoken commands would move closer to the broader goal of natural human-robot cooperation while still remaining grounded in explicit task semantics.

Integration with existing robot-programming environments is another direction. One possibility is to embed spatial authoring more directly into a teach-pendant or controller workflow, so that the same interaction becomes available as part of a conventional robot-programming process rather than as a separate frontend. This could make direct spatial input easier to integrate with established industrial workflows, without requiring it to exist as a separate frontend.

Formal evaluation is equally important. A natural continuation would therefore be a formal user study, together with stronger quantitative evaluation of tracking, scene estimation, and authoring performance. This would make it possible to answer more rigorously when spatial authoring reduces programming effort, for which task classes it is most useful, and where its practical limits remain.

Finally, the system should be exercised on a broader set of use cases. The present prototype already indicates that different task classes can benefit from the same architectural backbone for different reasons. Extending this work toward richer welding scenarios, dispensing and surface-processing tasks, or other human-in-the-loop workflows would help clarify the actual scope of the approach and would provide a stronger basis for both technical refinement and future evaluation.
